\chapter*{Notation}
\addcontentsline{toc}{chapter}{Notation}

\begin{ppdefn}
\ppterm{Speedup} is $S = T_1 / T_p$, where $T_1$ is serial time and $T_p$ is parallel time.
\end{ppdefn}

\begin{ppkey}
Always separate \ppterm{algorithmic parallelism} from \ppterm{implementation overheads}.
\end{ppkey}

\begin{ppdefn}
\ppterm{Vertex Set ($V$)}:
The set of all vertices (nodes) in a graph.
\end{ppdefn}

\begin{ppdefn}
\ppterm{Edge Set ($E$)}:
The set of all edges in a graph. Each edge may be represented as $(u,v)$ or $(u,v,w)$ for weighted graphs.
\end{ppdefn}

\begin{ppdefn}
\ppterm{Degree ($\deg(u)$)}:
The number of neighbours of a vertex $u$.
\end{ppdefn}

\begin{ppdefn}
\ppterm{Frontier}:
The set of vertices at the current BFS level being processed in parallel.
\end{ppdefn}

\begin{ppdefn}
\ppterm{Visited Array}:
A boolean array where $\texttt{visited}[v]$ indicates whether vertex $v$ has been discovered.
\end{ppdefn}

\begin{ppdefn}
\ppterm{Distance ($d[v]$)}:
The shortest-path distance from the source vertex to vertex $v$ in BFS.
\end{ppdefn}

\begin{ppdefn}
\ppterm{Compare-and-Swap (CAS)}:
An atomic operation $\textsf{CAS}(\&x, old, new)$ that updates $x$ to $new$ only if its current value equals $old$.
\end{ppdefn}

\begin{ppdefn}
\ppterm{Work ($W$)}:
Total number of operations performed by a parallel algorithm.
\end{ppdefn}

\begin{ppdefn}
\ppterm{Span / Critical Path ($T_\infty$)}:
Length of the longest chain of dependent operations in a parallel computation; equivalently, the running time on infinitely many processors.
\end{ppdefn}

\begin{ppdefn}
\ppterm{Level Work ($W_d$)}:
Total work at BFS level $d$, defined as
\[
W_d = \sum_{u \in \text{frontier}_d} \deg(u).
\]
\end{ppdefn}

\begin{ppdefn}
\ppterm{BFS Levels ($k$)}:
The number of levels (rounds) in a level-synchronous BFS from a chosen source --- equivalently, the eccentricity of the source.
\end{ppdefn}

\begin{ppdefn}
\ppterm{Graph Diameter}:
The maximum shortest-path distance over all pairs of vertices, $\max_{u,v} d(u,v)$.  Bounds the number of BFS levels for any source.
\end{ppdefn}

\begin{ppdefn}
\ppterm{DFS Tree}:
A spanning tree produced by a depth-first traversal; every non-tree edge of $G$ joins a vertex to one of its ancestors or descendants in the tree.
\end{ppdefn}

\begin{ppdefn}
\ppterm{Spanning Tree}:
A subgraph $T \subseteq G$ that is a tree and includes every vertex of $G$.  A \emph{minimum spanning tree} (MST) minimises the sum of edge weights.
\end{ppdefn}

\begin{ppdefn}
\ppterm{Cross-Component Edge}:
An edge $(u,v)$ such that $u$ and $v$ lie in different components, i.e., $\textsc{find}(u) \neq \textsc{find}(v)$.
\end{ppdefn}
\begin{ppdefn}
\ppterm{Component}:
A connected subgraph treated as a unit during Bor\r{u}vka rounds.
\end{ppdefn}

\begin{ppdefn}
\ppterm{Union--Find Structure}:
A data structure that maintains disjoint sets with operations \textsc{find} and \textsc{union}.
\end{ppdefn}

\begin{ppdefn}
\ppterm{Pointer Jumping}:
A parallel technique where
\[
\texttt{parent}[v] \leftarrow \texttt{parent}[\texttt{parent}[v]]
\]
is applied repeatedly to compress paths in union--find.
\end{ppdefn}

\begin{ppdefn}
\ppterm{Bor\r{u}vka Round}:
An iteration where each component selects its minimum-weight outgoing edge.
\end{ppdefn}

\begin{ppdefn}
\ppterm{Parallel Time ($T_p$)}:
Running time of an algorithm on $p$ processors.
\end{ppdefn}
